% Замените поясняющий текст собственными результатами.
% В шаблоне нет готового решения и обязательных внешних изображений.
\reportelement{Введение}

Сформулируйте цель и задачи лабораторной работы. Перечислите используемые
инструменты; персональные параметры укажите в таблице~\ref{tab:parameters},
изменив набор строк под своё задание.

\begin{table}[H]
  \caption{Персональные параметры}\label{tab:parameters}
  \centering
  \begin{tabularx}{\textwidth}{|p{0.46\textwidth}|X|}
    \hline
    Параметр & Значение\\ \hline
    Номер ИСУ & \StudentISU\\ \hline
    Персональный порт & 5455\\ \hline
    Персональный префикс & manuk\\ \hline
    Окружение & -\\ \hline
  \end{tabularx}
\end{table}

Результат хранится в ветке \texttt{lab-\LabNumber} своего форка:
\url{\RepositoryURL}. Замените адрес на действующий и сохраните коммит
версии сдачи.

\chapter{Ход выполнения}
\section{Подготовка окружения}

Опишите использованные программы и их версии. Ниже показано оформление
команд, а не готовый результат; документация Docker приведена как пример
источника~\cite{docker-docs}.

\begin{code}
# Замените на команды и результаты своего этапа.
docker version
docker ps -a
\end{code}

\section{Выполнение и проверка этапов}

Для каждого этапа приведите действие, наблюдаемый результат и объяснение.
На рисунке~\ref{fig:proof} показано место для собственного скриншота;
рамка не является доказательством выполнения работы.

% Добавьте свой файл step-01.png рядом с main.tex или замените путь.
% Пока файла нет, команда выводит рамку, но сборка продолжает работать.
\reportscreenshot{step-01.png}{Результат проверки этапа}{fig:proof}

Добавьте разделы по заданию своей лабораторной. Для ЛР №1 включите
таблицу состояний контейнера и остальные обязательные артефакты;
не считайте перечень примеров в этом макете полным заданием.

\section{Ошибки и анализ результатов}

Если возникала ошибка, опишите причину, выполненную проверку и исправление.
Не публикуйте настоящие пароли, токены и приватные ключи; в скриншотах
оставляйте требуемые заданием логин, хост и время, скрывая реальные секреты.

\reportelement{Заключение}

Сопоставьте результат с целью работы: что проверено, что получилось,
какие ограничения остались. Замените поясняющий текст собственными выводами.

% Оставьте только использованные источники и цитируйте их в тексте.
\begin{thebibliography}{9}
  \bibitem{docker-docs}
  Docker Documentation: документация Docker.
  URL: \url{https://docs.docker.com/}
  (дата обращения: укажите свою дату).
\end{thebibliography}

% Приложение необязательно. Удалите блок, если он не нужен для работы.
\startreportappendices
\reportappendix{Дополнительные команды и конфигурации}

Включите объёмные листинги, на которые ссылаетесь в основной части.
Короткий пример ниже показывает оформление, но не заменяет ваш код.

\begin{code}
# Свой безопасный пример конфигурации:
APP_SECRET=REDACTED
\end{code}
